\documentclass[10pt]{article}

\usepackage[margin=1in]{geometry}
\usepackage{amsmath,amssymb}
\usepackage{booktabs}
\usepackage{longtable}
\usepackage{array}
\usepackage{microtype}
\usepackage[shortlabels]{enumitem}
\setlist{nosep,leftmargin=*}
\usepackage{titlesec}
\titlespacing*{\section}{0pt}{1.1em}{0.5em}
\titlespacing*{\subsection}{0pt}{0.9em}{0.4em}
\titlespacing*{\paragraph}{0pt}{0.6em}{1em}
\usepackage[hidelinks]{hyperref}
\urlstyle{same}

\emergencystretch=1.5em
\newcommand{\decl}[1]{\texttt{#1}}
\newcommand{\Cmax}{C_{\max}}

\title{zk Payment Channels: a Definition, and the First\\
Machine-Checked Unlinkability Results for a Credit Construction}
\author{dmarz\thanks{Repository: \url{https://github.com/dmarzzz/zk-payments-confetti}. The definitions in this paper were produced and stress-tested under an agent-assisted workflow whose review protocol and full gate record are documented in the repository README.}}
\date{July 2026 --- draft}

\begin{document}
\maketitle

\begin{abstract}
A zk payment channel is a two-party channel in which the payee learns
nothing about payer identity across spends --- unlinkability within the
channel population --- while retaining balance security and double-spend
resistance. The object was named by BOLT in 2016, after which the line went
dormant: zkChannels stayed a draft, its implementation was archived in
2023, and no dedicated literature formed. Meanwhile the same shape has been
reinvented at least twice, as keyed-verification credit tokens (ACT, ARC)
and as ZK API Usage Credits. This paper defines the object as a tuple of
algorithms $(\mathsf{Setup}, \mathsf{Open}, \mathsf{Spend},
\mathsf{Redeem}, \mathsf{Close}, \mathsf{Dispute})$ with seven security
games, places it against the modern lineage, and reports a Lean~4
formalization of a flat-ticket RLN credit instantiation. Spend
unlinkability is machine-checked: against an adversarial payee holding
BOLT \S1.4's abort/evict powers, the advantage of distinguishing which of
two members emitted a whole epoch session is exactly zero --- to our
knowledge the first machine-checked unlinkability result for any
payment-channel or credit construction. No-overspend, both balance-security
theorems, closure liveness, a priced-divergence bound for an
eventually-consistent multi-gateway deployment, the exculpability (framing)
bound (under a stated random-oracle good-event hypothesis, with a
kernel-checked query-budget composition endpoint and the final
handler-coupling step still open), and the refund variant's safety and
conservation theorems --- now including the full failed-upgrade cascade and
finite-fleet aggregation --- are machine-checked with zero \texttt{sorry}
and no axiom beyond Lean's three standard ones. The model-to-real bridges
are no longer stated obligations: proof-bearing wire instances (a
masked-proof encoding, an interactive Sigma protocol, and a
lazy-random-oracle Fiat--Shamir compilation with explicit programming and
fork-collision bounds) each discharge the zero-knowledge bridge with zero
loss; an executable ledger refines every relational transition; a
multi-recipient portable-deposit accounting layer with a threshold-issuance
reference construction closes the definitional half of the named open
problem; and one-trace composition theorems deliver the channel and wire
guarantee bundles simultaneously on a single reachable trace. A third,
post-quantum instantiation --- Buterin's unidirectional nullifier-chain
channel --- is placed in the design space and its core is machine-checked
too: balance safety, both directions of the collision-based stale-close
mechanism, refund liveness, and per-request anonymity at advantage exactly
zero. The definition went through eleven rounds of
independent adversarial review, producing concrete counterexamples against
ten successive revisions; the repairs those
counterexamples forced --- gateway-bound messages, close-time netting,
verifiable spend counts at close --- are, we argue, design requirements for
any construction of this shape, and we present them as such.
\end{abstract}

\section{Introduction}

BOLT~\cite{greenmiers2017bolt} constructed anonymous payment channels: a
customer pays a merchant repeatedly off-chain, and the merchant learns
``no information beyond the fact that a valid payment \ldots\ has occurred
on a channel that is open with them.'' Then very little happened. BOLT
never shipped on Zcash~\cite{eccbolt2016}; its successor
zkChannels~\cite{zkchannels2021} never left DRAFT status; the
implementation, libzkchannels~\cite{libzkchannels}, was archived in
February 2023 as a proof of concept. There is no systematization of the
object, no formal definition of it as a distinct primitive, and --- as far
as a systematic search across TLA+, Why3, Rocq, Isabelle, and Lean can
establish --- no machine-checked proof of payment unlinkability for any
channel or credit construction in any prover. Every privacy proof in this
literature is pen-and-paper.

The object did not stay dead; it stayed unnamed. Keyed-verification credit
tokens --- ACT~\cite{act2025} and ARC~\cite{arc2025}, both 2025 IETF
drafts --- are anonymous balances spendable at exactly one verifier, with
double-spend prevented against the issuer's nullifier database: a
channel's state-update loop with the escrow replaced by ``the issuer was
already paid.'' ZK API Usage Credits~\cite{crapisbuterin2026credits} is a
deposit, a monotone spend index, and provider-countersigned balance
updates --- a commenter on the thread makes the channel mapping nearly
verbatim. These systems answer to the same security questions BOLT
answered to, but there is no shared definition to check them against.

The absence matters because this field has already demonstrated what
happens without one. A2L's privacy model~\cite{a2l2021} passed S\&P 2021
peer review; a year later Glaeser et al.~\cite{glaeser2022bcs} exhibited
two encryption schemes that satisfy its definitions and yield ``a
completely insecure system.'' The definitions, not the proofs, are the
risk surface.

\paragraph{What this paper does.} Three things.
\begin{enumerate}
\item \textbf{Defines the object} (\S\ref{sec:object}): a zk payment
channel is a tuple $(\mathsf{Setup}, \mathsf{Open}, \mathsf{Spend},
\mathsf{Redeem}, \mathsf{Close}, \mathsf{Dispute})$ over payers, payees,
and an idealized ledger, together with seven security games --- no
overspend, balance security for each side, spend unlinkability with an
abort/evict adversary, closure liveness, priced divergence, and
exculpability under collusion. The unlinkability game gives the
adversarial payee BOLT \S1.4's abort and eviction powers, because a game
without them proves a guarantee a real counterparty walks around. The
definition is scheme-agnostic; this is the contribution, and the rest of
the paper exists to defend it.
\item \textbf{Places the object} (\S\ref{sec:placement}) against the
modern lineage --- credit tokens, the hub-privacy line, ecash --- with
BOLT as origin, stating per system which property of the definition it
lacks.
\item \textbf{Formalizes and machine-checks} (\S\ref{sec:constructions}--\ref{sec:results}):
two instantiations --- a flat-ticket RLN credit protocol, arguably the
smallest machine-checkable unlinkability target in the literature, and a
refund-bearing variant covering variable-cost metering --- with the
safety, liveness, fleet, and \emph{privacy} theorems for the flat-ticket
instantiation kernel-checked in Lean~4 over VCV-io~\cite{vcvio2026}: the
spend-unlinkability advantage is proved exactly zero and the framing
probability at most $1/|F|$, alongside a built-in calibration pair that
separates the broken and fixed refund designs.
\S\ref{sec:results} states exactly what is proved and what is not.
\end{enumerate}

One methodological remark belongs up front. The definition in
\S\ref{sec:object} is revision eleven. Eleven independent review rounds
against earlier revisions produced concrete counterexamples against ten
successive revisions --- protocols and adversary schedules, not quibbles
--- and three of the
resulting repairs (gateway-bound messages, close-time netting for the
refund variant, verifiable spend counts at close) change protocol
behaviour, not merely its description. We present these as design
requirements discovered by adversarial definition review, with the
counterexamples that forced them, because any construction of this shape
that omits them is broken in the demonstrated way. The full review record,
with every counterexample, is in the repository~\cite{zkpcrepo}.

\paragraph{Honest scope.} The formalization covers the protocol layer over
an idealized ledger and idealized cryptography, in the random-oracle
model. It does not verify circuits, and \S\ref{sec:limits} lists the leaks
the theorems do not cover --- including two, the spend-count-at-close side
channel and within-epoch linkability, that are intrinsic to the design.

\section{The object}\label{sec:object}

This section mirrors the specification of record (\texttt{Spec.md}
revision~11 in the repository~\cite{zkpcrepo}); the Lean definitions are
traceable to it sentence by sentence, and it --- not the Lean --- is what
a reviewer reads.

\subsection{Setting and notation}

Three kinds of principal: \textbf{payers} (members), \textbf{payees}
(gateways --- $N$ of them may share one logical payee role), and an
\textbf{idealized ledger} $\mathcal{L}$: a totally ordered, atomically
executed transaction log that includes honest transactions within delay
$\Delta$ and runs the contract logic (escrow pool, membership tree, dedup
set, windows, slashing, automatic settlement) exactly as specified. Public
parameters: flat price $C$ (or per-spend cap $\Cmax$), deposit $D$, epoch
length $T_e$, per-gateway per-epoch budget $b$, fleet size $N$, end-to-end
reconciliation lag $L$, dispute window $\tau > \Delta$. Monetary
quantities and indices are naturals; the signal algebra lives in a prime
field $F_p$.

A member's long-term secret is $k \in F_p$ with identity commitment
$cm = H_{id}(k)$. The RLN signal at spend index $i$ on message $m$
is~\cite{rlndocs,crapisbuterin2026credits}
\[
a = H_a(k, i), \qquad x = H_x(m), \qquad y = k + a \cdot x, \qquad
nf = H_{nf}(a),
\]
with all hashes domain-separated and $H_x$ mapping into
$F_p \setminus \{0\}$ --- at $x = 0$ the signal is $y = k$, the secret
outright, a degenerate case the formalization isolated and the deployment
must exclude. Two well-formed signals on the same $(k, i)$ with
$x \neq x'$ are two points on a line: anyone computes
$a = (y-y')/(x-x')$ and $k = y - a \cdot x$. One signal, at $x \neq 0$, is
consistent with every candidate secret via a unique coefficient and
reveals nothing. This asymmetry --- machine-checked as
\decl{rln\_recover\_k} and \decl{rln\_single\_point\_hiding} --- is the
entire double-spend mechanism: reuse of an index is punished by public
recovery of $k$ and loss of the deposit, and the punishment is claimable
by anyone, with no watchtower. Tickets additionally carry an epoch
pseudonym $nf_e = H_e(k, e)$, linkable within an epoch by design (it is
what rate limiting counts) and unlinkable across epochs.

\textbf{Messages are gateway-bound.} A spend message is a pair
$m = (G, \hat{m})$ --- serving gateway identity plus payload --- and
$\mathsf{Redeem}$ at $G$ rejects tickets naming any other gateway. This is
a repair, not a transcription: see \S\ref{sec:repairs}.

\subsection{The algorithms}

\textbf{Setup} $\to pp$: CRS, hash descriptions, constants, the gateway
roster (sweep-authorized payee identities); for the refund variant, the
payee's signature keypair and a re-randomizable encryption key whose
decryption key is output to nobody.

\textbf{Open}$(pp, D) \to (cm, st_P)$: the payer samples $k$, posts
$(cm, D)$; the ledger appends $cm$ to the membership Merkle tree and adds
$D$ to a \emph{commingled escrow pool}. The pool is forced: before any
slash, a nullifier is by design unattributable to any $cm$, so no
per-channel draw is possible. Open is a public ledger event; hiding
\emph{that} a party opened a channel is out of scope (\S\ref{sec:limits}).

\textbf{Spend}$(pp, st_P, m) \to (t, st_P')$: off-ledger; the payer emits
a ticket $t = (\pi, root, e, nf_e, (x, y, nf))$ where $\pi$ proves, in
zero knowledge: membership of $cm$ under $root$; \emph{solvency}
$(i+1) \cdot C \le D$ at its current index $i$; and well-formedness of the
signal and epoch pseudonym for $(k, i, m, e)$. Emission consumes the
index. The honest retry rule after an abort is to re-send the
bit-identical ticket --- the same point on the same line, unslashable;
switching messages at a used index is self-slashing.

\textbf{Redeem}$(pp, st_G, t)$: the payee verifies $\pi$, the current
root, the epoch, the gateway binding, and the rate budget for $nf_e$; then
nullifier logic against its spent set: fresh $nf$ --- accept;
bit-identical tuple --- reject-duplicate (the abort-retry path); same
$nf$, different $x$ --- \emph{evidence}, forwarded to $\mathsf{Dispute}$.
In the fleet, gateways exchange accepted tuples within lag $L$, and a
gateway that merges a conflicting tuple emits evidence at merge time ---
required behaviour; without it a cross-gateway double spend that never
produces a third signal is never slashed.

\textbf{Close}: either side, any time. The payer's close and the payee's
settlement differ per instantiation and carry the two deepest repairs
(\S\ref{sec:repairs}, \S\ref{sec:constructions}). The payee side in the
flat instantiation is a unilateral \emph{sweep}: registered gateways
submit redeemed tuples; the ledger dedups by nullifier and pays $C$ per
fresh one from the pool.

\textbf{Dispute}$(pp, ev)$: anyone holding $ev = (nf, (x,y), (x',y'))$
with $x \neq x'$ submits it; the ledger recovers $a$ and $k$, checks
$nf = H_{nf}(a)$ and $cm = H_{id}(k)$ for an open channel, freezes the
channel, evicts $cm$ from the tree (the root rotates, so the member's
future spend proofs fail fleet-wide), and opens a claims window in which
registered gateways recover outstanding value --- sweeps first, then
documented conflicting acceptances, pro-rata within class --- against the
member's remaining deposit, \emph{but only for acceptances committed to
the ledger before the slash}. Each gateway checkpoints a binding Merkle
commitment to its accepted set at its own cadence (at least once per
epoch); claims open a pre-slash checkpoint with a membership witness. The
remainder goes to the evidence submitter as bounty. The
checkpoint-provenance rule exists because after a slash $k$ is public and
anything not anchored to a pre-slash commitment could be minted by the
fleet itself.

The specification distinguishes two kinds of slash, and the distinction is
load-bearing for the privacy analysis (\S\ref{sec:limits}). An
\emph{identity-slash} is the one just described: triggered by a
$\mathsf{Dispute}$ evidence pair, the ledger recovers $k$ from the line
algebra, $k$ becomes public, and the full MC4 claims window runs. A
\emph{fund-slash} --- the void-and-slash of a false unused-claim at close,
a settlement-detected sweep-bar violation, or a B failed-upgrade forfeit
--- carries no evidence pair, so $k$ stays hidden: the channel's remainder
stays in the pool with no submitter bounty (a bounty would need $k$ to
enumerate the member's other nullifiers), and ordinary sweeps continue.
Only the identity-slash publishes $k$; the close-dispute and settlement
slashes do not.

\subsection{The security games}

Seven statements, stated over the algorithms with explicit quantifiers,
adversary classes, and anti-vacuity witnesses; proof order
T1 $\to$ exculpability lemma $\to$ T2/T3 $\to$ T5 $\to$ T6 $\to$ T4 $\to$
T7. We summarize; the specification gives each in full.

\paragraph{T1 --- No overspend.} Against an arbitrary PPT payer coalition
controlling all payers and the scheduler, facing one honest payee (or a
perfectly synchronized fleet, $L = 0$): the accepted value attributed to
any secret $k$ never exceeds $D$. Attribution is by the
knowledge-soundness extractor. With lag $L > 0$ the invariant is genuinely
false during the window --- deliberately; the exact price is T6.

\paragraph{T2 --- Payee balance security.} An honest payee following the
sweep protocol (including monitoring $\mathsf{Dispute}$ windows, and
keeping its checkpoints current at payer-closes) settles exactly
$C \cdot |T|$ for its accepted set $T$, against payers who close early,
close with false unused-claims, or deliberately trigger their own slash.
The upper bound is unconditional (ledger dedup plus per-ticket price).

\paragraph{T3 --- Payer balance security.} Against a malicious payee
coalition (all $N$ gateways) plus all other payers: an honest payer that
emitted $j$ tickets recovers at least $D - jC$ (flat;
$D - j\Cmax + R$ with held receipts in the refund variant), and no
adversary produces $\mathsf{Dispute}$ evidence that slashes an honest
payer except with negligible probability. Emission is the authorization
event: a payee that takes a ticket and refuses service is paid for it, and
the theorem states that abort-griefing cost explicitly rather than hiding
it.

\paragraph{T4 --- Spend unlinkability (the headline game).} The adversary
plays the payee --- all $N$ gateways, all payee keys, arbitrarily many
corrupt payers. The challenger runs two honest candidates with equal
deposits, opened in batch. Pre-challenge, the adversary drives both
candidates freely: spends on messages of its choice, retries, receipt
issuance (refund variant), and the \emph{abort/evict powers} of BOLT
\S1.4~\cite{greenmiers2017bolt} --- it may refuse service to a candidate
from any point on, and in the refund variant withholding receipts can
drive a candidate insolvent. At challenge time the adversary outputs a
message \emph{vector} $\vec{m}^* = (m^*_1, \ldots, m^*_q)$ of its choice
($q \ge 1$); the game checks the current epoch $e^*$ is fresh for both
candidates and both are challenge-capable for $q$ (open, unclosed, solvent
for $q$ more spends); a candidate evicted into insolvency makes the
challenge return $\bot$ --- the game charges eviction to the anonymity set,
not the scheme, which is precisely the calibrated content of the abort
attack. Otherwise the hidden candidate $P_b$ emits the whole
$q$-spend session $t^*_1, \ldots, t^*_q$ inside $e^*$ (sharing the one
session pseudonym $nf_{e^*}$), the adversary receives the batch,
\emph{and the game ends}: the guess is a pure function of retained memory
plus the challenge response. Advantage is $|\Pr[b' = b] - \tfrac12|$, with
$b$ sampled at game start so $\bot$-paths contribute exactly $\tfrac12$.
The session form --- an adversary-chosen vector rather than a single
spend --- is itself a repair the external review forced: the single-spend
game certifies only unlinkability of a member's \emph{first} spend of an
epoch, and a scheme leaking a persistent cross-epoch tag on second-and-
later spends would pass it while being lifetime-linkable for any member
that spends twice in an epoch (the deployed fleet's normal usage). The
session challenge exercises the second-spend wire format and catches that
class; what is certified is the unlinkability of a member's \emph{whole
epoch session} to its identity and its other epochs, within-session
linkage via $nf_{e^*}$ remaining by design.

Two further features of this game are results of review, not first drafts.
First, the natural game with post-challenge oracles is
\emph{unsatisfiable}:
three distinguishers win it against every scheme including sound ones ---
replay the challenge through the retry oracle, probe which candidate
exhausts solvency one index early, read the spend count off a later
close. All three exploit the bit-dependent continuation, so the game
terminates at challenge delivery, and the residual leak (aggregate spend
count at close) is stated honestly in \S\ref{sec:limits} rather than
defined away. Second, the game carries a \emph{calibration requirement}:
instantiated on the refund variant with a static encrypted running total
(the original design), a concrete distinguisher must win --- it issues
receipts with equal totals to both candidates and bit-matches the
presented ciphertext against its own issuance
transcript~\cite{omarespejel2026} --- and instantiated on the
re-randomized repair, all PPT advantage must be negligible. A game that
cannot separate the broken variant from the fixed one is the wrong game;
this pair is the built-in test that ours is not, and the broken direction
must be a constructive term in the formalization, not an unproven gap.

\paragraph{T5 --- Closure liveness.} An honest closer settles by
$t + \Delta + \tau$ exactly, against a counterparty that never appears;
silence, garbage, and concurrent disputes cannot extend the bound.

\paragraph{T6 --- Priced divergence (fleet).} Against one corrupted member
facing $N$ honest gateways with end-to-end reconciliation lag $L$: total
accepted value is at most
\[
\lfloor D/C \rfloor \cdot C + f(L), \qquad
f(L) = N \cdot b \cdot (\lceil L/T_e \rceil + 1) \cdot C,
\]
and a member with two conflicting accepted signals is slashed fleet-wide
within $L$ of the second acceptance. The deployment condition $f(L) < D$
bounds the maximal burst by one deposit. The theorem deliberately does
\emph{not} claim attacker unprofitability or universal recovery: recovery
through the claims window is capped by the remaining deposit and gated on
pre-slash checkpoints, and an exhaust-then-burst schedule leaves a
near-empty remainder. The operational levers are sweep cadence and
checkpoint cadence. This is the theorem that makes an
eventually-consistent spent set safe to run: the async window is priced by
the deposit, not trusted away.

\paragraph{T7 --- Exculpability under collusion (fleet).} $N-1$ gateways
pooling all transcripts, plus corrupt members, with adaptive spend and
close oracles against one honest member, produce evidence that slashes it
with probability at most negligible. The algebraic core: the coalition
holds at most one point per line, and one point at $x \neq 0$ determines
nothing about $k$; producing a second point is computing $k$. This is what
makes an automatic, anyone-can-submit slash safe against the protocol's
actual threat model --- the fleet's own operators.

\subsection{Three repairs, with the counterexamples that forced
them}\label{sec:repairs}

The specification records twenty modeling choices; three changed protocol
behaviour and generalize beyond this protocol. Each was found by a
reviewer who did not write the definition, as a concrete adversary
schedule against the previous revision.

\paragraph{Replay across payees (gateway-bound messages).} In the first
revision, a spend message was just the request payload. The
counterexample: replay the bit-identical ticket at all $N$ gateways. Same
$x$, same $y$, same nullifier --- no conflicting pair ever forms, no
evidence exists, the slash clock never starts, and each gateway's private
spent set happily accepts its copy. Excess extraction is roughly
$(N-1) \cdot D$ against a claimed bound of order $r \cdot L \cdot C$
(numerical witness from the review record: $N{=}3$, $b{=}100$,
$T_e{=}1$d, $L{=}1$\,min, $D{=}1000C$ gives $2000C$ of excess against a
bound of $0.2C$). The repair binds the serving gateway's identity into the
hashed message, so cross-gateway reuse of an index \emph{forces} distinct
$x$ --- a conflicting pair --- and the detection argument becomes sound.
The general lesson: in any multi-verifier deployment of a nullifier
scheme, replay across verifiers is value-bearing unless the verifier
identity is inside the proof's message. Three independent reviewers found
this; it is the record's canonical wrong-definition-nearly-proved case.

\paragraph{The gap-index close (verifiable spend counts).} Through five
revisions, a payer closed by declaring its spend count, with
understatement policed by collision with its own spent indices. The
fifth-round counterexample: nothing enforces index \emph{contiguity} ---
indices are hidden witnesses --- so a payer skips index 0, spends at
indices $1..m$, and closes declaring index 0. The declaration collides
with nothing, is undisputable, and recovers the full deposit after
consuming the service. The root cause is structural: the ledger has no
verifiable spend count. The repair is different per instantiation, and the
difference is itself a finding (\S\ref{sec:constructions}): the flat
instantiation closes by \emph{enumerating} PRF-fresh nullifiers of
claimed-unused indices (false claims are disproven by bit-match against a
pre-close checkpoint; genuinely unused nullifiers are PRF-hidden before
the close reveals them, so no pre-close checkpoint can contain them ---
the same mechanism protects honest closers from fabricated disputes); the
refund instantiation certifies the count inside the receipt chain and
proves each spend's index equals its receipt's count, making spends
contiguous by construction, then reveals the nullifier of the first index
\emph{beyond} the declared count so that closing on a stale receipt
self-convicts.

\paragraph{Fund conservation under refunds (close-time netting).} The
flat instantiation lets gateways sweep per accepted nullifier. The
second-round counterexample proved this cannot coexist with refunds:
sweeps pay $\Cmax$ per nullifier \emph{and} the close pays the refund
total $R$ --- $D + R$ out of a $D$ deposit --- and pre-slash
unattributability makes per-channel netting impossible, so payee balance
security, payer balance security, and pool solvency are jointly
unsatisfiable. The repair removes unilateral sweeps from the refund
variant entirely: each channel settles exactly once, at close, with the
payer receiving $(D+R) - j \cdot \Cmax$ and the payee
$j \cdot \Cmax - R = \sum_\ell c_\ell$, under two circuit-enforced caps
($R \le j \cdot \Cmax$, else a colluding payee signs inflated refunds and
drains the pool; $j \cdot \Cmax \le D + R$, else an overstated count pays
the payee unboundedly). Conservation is then exact by construction. The
general lesson: refund privacy and unilateral per-ticket settlement are
incompatible under a commingled escrow; a refund-bearing anonymous channel
must settle per channel, at close.

\subsection{Model boundary}\label{sec:boundary}

We formalize the protocol layer over the idealized ledger and idealized
cryptography, in the random-oracle model. The NIZK relation is the
mathematical statement above; the fidelity of any circuit to it is out of
scope --- the repository's own wording is that anyone claiming it verifies
SNARKs is misreading it. Named assumptions (knowledge soundness, zero
knowledge, PRF/ROM with domain separation, EUF-CMA receipts,
re-randomizable encryption, and a separately named
\decl{single\_signal\_hiding} --- the additive, KDM-flavoured key use in
$y = k + a \cdot x$ is not implied by standard PRF security) are confined
to one audited registry file. Explicitly out of scope: circuit
correctness, network-level timing and content fingerprints, funding-graph
leakage at $\mathsf{Open}$, a global passive adversary, relationship
anonymity (the transport layer's property), and the policy stake of the
source construction.

\section{Placement}\label{sec:placement}

Table~\ref{tab:placement} compares against the modern lineage, with BOLT
as origin and --- via its \S1.4 abort attacks --- the sharpest
transferable threat model. ``Ours'' means the object of
\S\ref{sec:object}: recipient-bound anonymous spending with balance
security both sides, detect-and-slash double-spend handling, and
unlinkability against a payee holding abort/evict powers. Rows are sourced
from a verified literature sweep (\texttt{RESEARCH.md} in the
repository~\cite{zkpcrepo}); each cites its primary document.

\begin{footnotesize}
\begin{longtable}{>{\raggedright\arraybackslash}p{0.26\textwidth}>{\raggedright\arraybackslash}p{0.32\textwidth}>{\raggedright\arraybackslash}p{0.34\textwidth}}
\caption{Placement against the lineage.}\label{tab:placement}\\
\toprule
System & What it is & What it lacks relative to the \S\ref{sec:object} object \\
\midrule
\endfirsthead
\toprule
System & What it is & What it lacks (cont.) \\
\midrule
\endhead
\bottomrule
\endfoot
Chaum blind signatures~\cite{chaum1983}; Chaum--Fiat--Naor~\cite{cfn1990} &
Bearer ecash; CFN adds offline double-spend detection revealing the cheater &
Custodial mint; no escrow or dispute game; no balance-security theorems against the verifier; detection reveals identity but recovery is out of band \\
\addlinespace
Compact E-Cash~\cite{compactecash2005} &
$O(1)$ wallets, serial-number nullifiers, identity extraction on reuse --- BOLT's direct ancestor &
Same custodial shape; no channel state, no close/dispute, no priced asynchrony \\
\addlinespace
BOLT~\cite{greenmiers2017bolt} &
Named the object: uni/bidirectional anonymous channels, revocation punishment, exculpability; \S1.4 states the abort attacks &
Punishment requires chain-watching (watchtowers); amounts visible to the merchant; anonymity conditional on anonymized funding; no machine-checked statement; abort attacks stated, not carried into the security games \\
\addlinespace
zkChannels~\cite{zkchannels2021} &
BOLT over an abstract arbiter, Pointcheval--Sanders signatures, an unlinking protocol for the open &
Draft; dispute forfeits the entire balance to the merchant (a griefing cliff our close/dispute avoids); both sides online-or-watchtowered; archived implementation~\cite{libzkchannels} \\
\addlinespace
TumbleBit~\cite{tumblebit2017} &
Unidirectional payment hub; the field's reference anonymity accounting ($k$ = completed payments per epoch) &
Hub topology --- a third party that is neither payer nor payee; fixed denominations; unlinkability vs the hub, not vs the payee \\
\addlinespace
A2L~\cite{a2l2021} / A2L+~\cite{glaeser2022bcs} &
Adaptor-signature hub; A2L+ repairs the 2021 model after definitional counterexamples &
Hub topology, hub adversary; the 2021 episode is this paper's cautionary tale rather than its competitor \\
\addlinespace
BlindHub~\cite{blindhub2023} &
Variable-amount hub unlinkability via blind adaptor signatures &
Hub topology; solves amount-as-linking-tag, which flat pricing avoids by construction \\
\addlinespace
Accio~\cite{accio2023} &
Hub unlinkability, variable amounts, no NIZKs, no pre-locking &
Hub topology; payer side plaintext; no payee-adversary unlinkability \\
\addlinespace
Adaptor-signature foundations~\cite{adaptor2024} &
Formal definitions for the hub line's main tool &
A primitive, not a channel; listed because it is where that line's definitional rigour now lives \\
\addlinespace
Cashu~\cite{cashu} / Fedimint~\cite{fedimint} / Taler~\cite{taler2016} &
Deployed Chaumian ecash: single mint / BFT federation / exchange with unlinkable change &
Online prevention needs a synchronous spent set (Fedimint pays consensus latency on the spend path); custodial; no dispute game, no payee-balance theorem against the mint \\
\addlinespace
Privacy Pass~\cite{rfc9576} &
Standardized single-use unlinkable tokens &
Unit tokens, not balances; replay defense still needs a per-origin spent list; no value semantics \\
\addlinespace
ARC~\cite{arc2025} &
Keyed-verification $N$-use credential, presentations pairwise unlinkable &
Issuer $=$ verifier with online prevention: no escrow, no dispute, no slash; nothing prices a stale verifier view (our T6 has no analogue) \\
\addlinespace
ACT~\cite{act2025} &
ARC with money semantics: hidden balance, nullifier per spend, blind change --- the closest living relative &
Prevention against the issuer's DB, so nothing to say about multiple verifiers, asynchrony, or framing; no exculpability game (the issuer's DB is trusted); no close/liveness \\
\addlinespace
Nym zk-nym~\cite{nymzknym} &
Threshold-issued ticketbooks (Coconut lineage~\cite{coconut2018}), spendable at any gateway, deferred Bloom-filter reconciliation &
Closest deployed shape to our fleet setting, but reconciliation cadence, synchrony assumptions, and double-spender penalties are undocumented --- exactly the parameters our T6 makes explicit and proves against \\
\addlinespace
ZK API Usage Credits~\cite{crapisbuterin2026credits} &
Deposit + monotone index + RLN detect-and-slash + refund tickets; the construction of record for our refund variant &
No formal definition or games; settlement cadence underspecified; the self-slash race and gap-index close (\S\ref{sec:repairs}) live in the gap between its prose and a full specification \\
\addlinespace
PrivateX402~\cite{privatex402} &
One deposit funding $N$ recipient channels; allocation privacy vs chain observers &
No unlinkability vs the recipient at all (stable session key, cumulative receipts) --- the property that defines our object \\
\addlinespace
Nirvana~\cite{nirvana2022} / RRTE~\cite{rrte2023} &
Anonymous zero-confirmation payment guarantees; double-spender revealed by threshold decryption &
Threshold committee in the loop; retail topology; no channel state or close \\
\end{longtable}
\end{footnotesize}

Two structural observations the table compresses. First, the hub line
(TumbleBit through Accio) solves a different problem: it hides
who-pays-whom from a \emph{third party}; in our setting the payee itself
is the adversary, so the right ancestor is BOLT, not the tumblers. Second,
the prevention/punishment split is the live design axis: ACT/ARC and the
ecash systems prevent double-spends against a synchronous verifier
database, which is exactly what a fleet of mutually distrusting verifiers
cannot cheaply have; BOLT punished via revocation, which needs
watchtowers; the RLN line makes punishment cryptographic-economic and
watchtower-free, at the price of a detection window that T6 prices.

On the verification side: Lightning has a UC
treatment~\cite{kiayiaslitos2020}, machine-checked fund safety in
Why3~\cite{why3ln2025}, and TLA+ model-checking made feasible by proven
refinements~\cite{tlaln2025} --- balance security is covered ground.
Privacy is not: the mature game frameworks are SSProve~\cite{ssprove2021}
(Rocq) and CryptHOL~\cite{crypthol2017} (Isabelle), and none of them has
been used on channel unlinkability. We work in Lean~4 over
VCV-io~\cite{vcvio2026}, which with ArkLib~\cite{arklib} is already
verifying SNARK components --- the ecosystem this protocol stack would
eventually rest on.

\section{The instantiations}\label{sec:constructions}

\paragraph{A: flat-ticket RLN credits.} Deposit $D$, flat price $C$,
solvency $(i+1) \cdot C \le D$, per-index nullifier, slash on reuse. No
refunds, no revocation, one inequality; $\mathsf{Spend}$ is a single
message and the payee holds no per-payer state, so a payee abort is
exactly a denial of service. This is the protocol a Tor egress
fleet~\cite{egresspost2026} runs --- near-uniform request cost makes the
refund machinery deletable --- and it is deliberately the smallest
instantiation with all seven games defined on it. The fleet theorems
T6/T7 exist only here: $N$ gateways with local spent sets, gossip within
$L$, merge-time evidence.

\paragraph{B: refund-bearing credits.} The variant the source
construction~\cite{crapisbuterin2026credits} needs for LLM-style cost
variance (the thread's argument: without refunds, a budget covering
worst-case cost overshoots the mean by orders of magnitude). Per accepted
spend the payee declares actual cost $c \le \Cmax$ and certifies a receipt
chain over the triple $(tag, R, n)$ --- chain tag binding receipts to the
channel (without it, receipts farmed on a cheap channel splice into
another channel's solvency proofs, a first-round counterexample), running
refund total $R$, and certified spend count $n$. Solvency becomes
$(i+1) \cdot \Cmax \le D + R$, proven against the certified ciphertext,
with the spend's index proven equal to $n$. Both representations of the
certified total are formalized: \textbf{B-static} (present the last-signed
ciphertext bit-identically --- the original design, broken: the payee
bit-matches it against its own issuance transcript~\cite{omarespejel2026},
and the genesis receipt even links the first spend to the identity) and
\textbf{B-rerand} (re-randomize and prove equivalence in zero knowledge
--- the patch). The unlinkability game \emph{does} fail on B-static (a
concrete distinguisher at advantage $\tfrac12$) and pass on B-rerand
(advantage $0$); this calibration pair is machine-checked
(\S\ref{sec:results}), the built-in evidence that the game is not vacuous.

\paragraph{The close asymmetry, as a design-space observation.} The
gap-index counterexample (\S\ref{sec:repairs}) hit both instantiations,
and they cannot share a repair, because their information structures
differ. A's spending is non-interactive --- one message, no
countersignature --- so no certified count can exist, and the close must
\emph{reveal}: enumerate the unused indices' nullifiers, sized
$\lfloor D/C \rfloor - j$, with disputes by checkpoint bit-match. B's
spending is interactive, so the count can be certified in the receipt
chain and the close is $O(1)$: present the latest receipt, reveal one
nullifier past the count. Non-interactive spending buys abort-immunity and
no per-payer payee state, and pays for it at close time with an
enumeration proportional to the unspent budget; interactive receipts buy
certified counts and cheap closes, and pay with a live abort lever
(withheld receipts stall solvency). We know of no prior statement of this
trade-off, presumably because no prior design was pushed through an
adversarial review that forced both closes to be sound.

\paragraph{C: the unidirectional nullifier-chain channel.} A third
instantiation, due to Vitalik Buterin (communicated via
dmarzzz\footnote{\url{https://gist.github.com/dmarzzz/ddcd1302c5f511001f8f46102874a08e};
archived in the repository as
\texttt{research\_knowledge/vitalik-nullifier-chain-channel.md}.}), sits
at a different point in the design space and is markedly simpler; we
record it here and formalize its core alongside A and B. Alice maintains a
nullifier chain $N_{i+1} = H(N_i, c)$ under a private seed $c$.
$\mathsf{Open}$ deposits $D$, names the recipient publicly, and commits
(hiding) to $N_1$. Each payment reveals the parent state's committed
next-nullifier and carries a ZK proof that the parent is either the
on-chain genesis or \emph{some} Bob-countersigned state --- the signature
verified inside the proof, so \emph{which} state stays hidden --- with
$\mathit{parent\_balance} + \delta = \mathit{new\_balance} \le D$,
balances hidden, $\delta$ public; the message commits to the new balance
and a fresh next nullifier, and Bob countersigns unless he has seen the
revealed nullifier before. $\mathsf{Close}$ opens the closed state's
committed next-nullifier on chain; Bob challenges by exhibiting a message
that revealed the same nullifier --- a collision proves the closed state
was extended --- and a successful challenge forfeits the whole deposit to
him. The one mechanism, the nullifier chain, does both jobs: duplicate
detection during payments and stale-close detection at close, uniformly
down to the genesis-refund case.

The comparison is instructive in both directions. The design is
unidirectional, per-pair, and interactive (every payment needs a
countersignature --- B's abort lever, not A's abort-immunity), and its
base form leaks $D$ and the final split at the channel boundaries;
recipient anonymity, deposit privacy, and close privacy are add-ons (an
ephemeral per-channel key, or shielded-pool-integrated opens and closes),
and it inherits the recipient-boundness of \S\ref{sec:limits} unchanged.
In exchange it deletes machinery this paper spent sections on: no epochs,
no rate-limiting algebra, no fleet, no refund cascade --- and, most
notably for \S\ref{sec:limits}, \emph{no identity-slash}. Its penalties
are fund-forfeit only, so the retroactive deanonymization that FRAME must
price in the RLN design (a published $k$ unlocks the member's entire
request history) has no analogue here: a cheating Alice loses money, not
her lifetime privacy, and the FRAME-shaped question collapses to the
collision soundness of the challenge --- an honest Alice's latest-state
close opens a nullifier no message ever revealed. The anonymity claim is
also differently scoped: per-request unlinkability toward the recipient
(Bob cannot link two payments to a sender or channel), where T4 is
population unlinkability toward a payee across members --- hidden balances
are what make it hold (a visible cumulative balance would let Bob rejoin
the chain from his own $\delta$ records, which high-entropy $\delta$ makes
easier, not harder). The construction is also natively post-quantum
(hashes and STARK-friendly signatures; no curves), which none of A/B's
reference wires are. Its Lean formalization (\texttt{Zkpc/Chain/}) follows
the same split this paper uses everywhere --- safety and stale-close
collision detection as Class-A invariants, per-request anonymity as a
Class-B coupling, with the STARK and signature scheme idealized at the
same boundary as A and B's proofs --- and is machine-checked
(\S\ref{sec:results}).


\section{Results: what is machine-checked}\label{sec:results}

Everything in this section refers to Lean declarations in the
repository~\cite{zkpcrepo}, building with zero \texttt{sorry} under CI
that also forbids \texttt{axiom} outside the assumptions registry and
forbids \texttt{admit}/\texttt{native\_decide}. The development in fact
contains \emph{no \texttt{axiom} declarations at all}: each named
assumption is discharged by construction in the idealized model (knowledge
soundness as transition guards --- an accepted ticket \emph{is} its
extracted witness; zero knowledge as proof-free adversary views; hashes as
lazily sampled random oracles), so \texttt{\#print axioms} on every
theorem shows only Lean's
\texttt{propext}/\texttt{Quot.sound}/\texttt{Classical.choice}. This was
audited declaration by declaration (K2): the unlinkability, framing,
calibration, and refund theorems each reduce to exactly those three
standard Lean axioms and nothing more. The trust surface is the
definitions, which is where the review effort went.

The model boundary bears repeating before the list: idealized ledger,
random-oracle model, no circuits. These are theorems about the protocol
layer.

\paragraph{Machine-checked theorems (flat-ticket instantiation).}
\begin{itemize}
\item \textbf{T1} --- \decl{T1\_no\_overspend}
  (\texttt{Zkpc/Core/T1.lean}): accepted value per secret never exceeds
  $D$, single honest payee, arbitrary payer coalition.
\item \textbf{Exculpability lemma, symbolic form} ---
  \decl{honest\_never\_slashed} (\texttt{Zkpc/Core/T1.lean}): a
  protocol-following payer is never slashed in any reachable state; the
  invariant carrier is \decl{reach\_inv}.
\item \textbf{T2} --- \texttt{Zkpc/Core/T2.lean}: unconditional upper
  bound \decl{T2\_upper} / \decl{T2\_paid\_exact} (ledger dedup plus
  per-ticket price cap the payee's revenue), collectability
  \decl{T2\_collectable} (from any reachable state a finite sweep sequence
  settles every accepted-unswept ticket whose window is open, with no
  payer cooperation), and \decl{T2\_settles\_exactly}.
\item \textbf{T3} --- \decl{T3\_settled\_amount},
  \decl{T3\_payer\_balance\_security} (\texttt{Zkpc/Core/T3.lean}): the
  settled amount satisfies $paid + j \cdot C = D$ exactly (stronger than
  the specification's floor), plus no-framing via the exculpability
  lemma.
\item \textbf{T5} --- \decl{T5\_payer\_close\_liveness}, with persistence
  \decl{settleClose\_stable} and progress \decl{tick\_progress}
  (\texttt{Zkpc/Core/T5.lean}): settlement is enabled from window expiry,
  stays enabled under every adversarial action, and a settling
  continuation exists from every reachable state; the payee half is T2's
  collectability.
\item \textbf{T6} --- \decl{T6\_priced\_divergence},
  \decl{T6\_accept\_count}, \decl{T6\_slash\_within\_L}
  (\texttt{Zkpc/Fleet/T6.lean}), on the counting lemmas
  \decl{card\_le\_solvency\_of\_conflictFree} and
  \decl{card\_le\_rate\_window} and the epoch-straddle lemma
  \decl{epochs\_in\_window} (\texttt{Zkpc/Fleet/Basic.lean}): accepted
  value $\le \lfloor D/C \rfloor \cdot C + N b (\lceil L/T_e \rceil + 1)
  C$ and fleet-wide slash within $L$. The formalization surfaced two
  boundary facts the prose missed: the ticket-count form is false at
  $C = 0$ (a zero-price fleet accepts unboundedly many non-conflicting
  tickets), and $T_e > 0$ is load-bearing.
\item \textbf{RLN algebra} --- \texttt{Zkpc/Games/RLN.lean}:
  \decl{rln\_recover\_a}, \decl{rln\_recover\_k} (two points reveal the
  secret --- $\mathsf{Dispute}$'s completeness),
  \decl{rln\_single\_point\_hiding} (one point at $x \neq 0$ is consistent
  with every candidate secret via a unique coefficient --- the
  information-theoretic core of \decl{single\_signal\_hiding}),
  \decl{rln\_evidence\_complete}, \decl{rln\_evidence\_sound} (the
  ledger's slash predicate is satisfiable only by a genuine two-point
  exposure), and \decl{rln\_x\_zero\_degenerate} (the $x = 0$
  counterexample that forced the domain-separation requirement).
\end{itemize}

\paragraph{Machine-checked privacy theorems (the headline).} The two
privacy games are defined over the advantage framework of
\texttt{Zkpc/Games/Framework.lean} (\decl{guessGap}, the
$|\Pr[b'=b] - \tfrac12|$ form; the challenge-terminated adversary shape
\decl{ChalAdversary}, whose post-challenge oracle silence holds by type;
the abort/evict wrapper \decl{withEvict}), and both are proved.
\begin{itemize}
\item \textbf{T4 --- spend unlinkability, exactly zero advantage} ---
  \decl{T4\_flat\_unlinkability} (\texttt{Zkpc/Games/T4.lean}), over the
  session-form game \decl{unlinkGame}/\decl{unlinkAdvantage}
  (\texttt{Zkpc/Games/Unlink.lean}): for the flat-ticket ideal
  instantiation, \emph{every} UNLINK adversary at \emph{every} deposit
  budget has advantage $= 0$ --- information-theoretically perfect
  unlinkability of a member's whole epoch session, against a payee with
  the full pre-challenge abort/evict power of the BOLT \S1.4 oracle. To
  our knowledge this is the first machine-checked spend-unlinkability
  result for any payment-channel or credit construction. The two honest
  residues (within-session $nf_{e^*}$ linkage, MC6; the spend count at
  close, MC15) are pinned, not hidden: \decl{flat\_closeViewSimulatable}
  proves the close view is simulatable from $(cm, \text{count})$ alone, so
  the leak is exactly the count and no more.
\item \textbf{T7 --- exculpability (framing) bound} ---
  \decl{T7\_frame\_bound} (\texttt{Zkpc/Games/T7.lean}), over
  \decl{frameGame}/\decl{frameWinProb} and the win predicate
  \decl{Slashes} (deliberately stronger than $\mathsf{Dispute}$): under
  the stated random-oracle good-event hypothesis \decl{hobliv} (the
  adversary's evidence is independent of the secret $k$), the probability
  that an $N-1$-gateway coalition frames an honest member is at most
  $1/|F|$ --- the blind-guess floor. The query-term accounting that would
  discharge \decl{hobliv} is now itself mostly kernel-checked
  infrastructure rather than prose: structural per-channel query budgets
  (\decl{FrameQueryBounds}), the adaptive uniform-secret first-hit bound
  \decl{uniformSecretProbeBound} ($q/|F|$), the adaptive multi-target
  slope-preimage bound \decl{uniformSlopeProbeBound}
  ($q_{Nf}\,q_{sig}/|F|$), a quantitative real-to-ideal assembly theorem
  \decl{T7\_frame\_bound\_of\_pointwise}, and the composition endpoint
  \decl{T7\_frame\_query\_bound}, which turns a deferred-sampling
  certificate into the corrected concrete bound
  $(q_A + q_E + q_{Id} + q_{Nf} q_{sig} + q_{sig}^2 + 1)/|F|$. An exact
  audited ghost of the real handler (\decl{FrameAudit}: bad-event
  monotonicity, per-step resource growth, projection to \decl{frameImpl}
  on every adaptive run) and a secret-independent ideal handler with a
  canonical secret-erasing state map and per-oracle idealization step
  lemmas (\decl{FrameIdeal}) are proved. The certificate statement itself
  was then stress-tested and \emph{refuted} as originally frozen
  (\decl{frameDeferredSampling\_refuted},
  \texttt{Zkpc/Games/FrameDeferred.lean}): a kernel-checked two-probe
  adversary --- two $H_{id}$ queries against the public commitment ---
  forces any single secret-independent generator to carry mass
  $\ge 1 - 2/|F|$ on each of two disjoint slash events, so the
  pointwise-in-$k$ domination demanded by \decl{FrameDeferredSampling} is
  unsatisfiable for $|F| > 5$. This is a gate finding in the repository's
  sense (the definition, not the theorem, was wrong): the leak mass is
  bounded only \emph{averaged over the uniform secret}, which is also
  exactly the quantity the FRAME experiment consumes. The corrected socket
  \decl{FrameDeferredSamplingAvg} composes to the \emph{same} final bound
  via \decl{T7\_frame\_query\_bound\_avg}, and the pointwise form strictly
  implies it. Honestly in-place: constructing the averaged certificate
  from \decl{frameImpl} by the stateful transcript coupling remains open,
  and we do not claim the unconditional bound --- but the pointwise route
  is now provably a dead end rather than an open problem, which is itself
  information. Three \emph{must-win} calibration adversaries confirm the
  game has teeth: \decl{frameWinProb\_YK\_eq\_one} (degenerate $y = k$),
  \decl{frameWinProb\_aReuse\_eq\_one} ($a$ reused across indices), and
  \decl{frameWinProb\_slopeReveal\_eq\_one} (the formal witness that the
  slope-preimage channel is real, so the $q_{Nf}\,q_{sig}$ term is
  required, not conservative slack) each frame with probability exactly
  $1$.
\item \textbf{The calibration pair (the built-in definitional test)} ---
  \decl{unlinkAdvantage\_staticDistinguisher\_eq\_half} (the broken
  B-static design: a concrete distinguisher wins at exactly $\tfrac12$)
  and \decl{unlinkAdvantage\_bRerand\_eq\_zero} (the B-rerand repair: every
  adversary at exactly $0$), in \texttt{Zkpc/Games/Calibration.lean}. The
  \emph{same} game separates the broken variant from the fixed one, which
  is the property a wrong game lacks. The must-catch battery joins it:
  \decl{unlinkAdvantage\_aIndexLeak} (index-in-clear, A's first
  calibration point), \decl{unlinkAdvantage\_nfeReuse} ($nf_e$ derived
  without $e$), and \decl{unlinkAdvantage\_multTagDistinguisher\_eq\_half}
  (the multiplicity-tag scheme the session form was introduced to catch,
  won by a $q = 2$ distinguisher) --- each at exactly $\tfrac12$. The
  algebraic core T7 reduces to (\decl{rln\_single\_point\_hiding},
  \decl{rln\_evidence\_sound}, \texttt{Zkpc/Games/RLN.lean}) is in the
  flat-ticket list above.
\end{itemize}

\paragraph{Machine-checked refund-variant safety, cascade, and fleet
aggregation.} The refund variant (instantiation~B) is machine-checked in
\texttt{Zkpc/Refund/}: \decl{T1\_B\_no\_overspend} (accepted cost
$\sum_\ell c_\ell \le D$), \decl{T3\_B\_floor} (a cooperatively-settled
payer recovers exactly $D - \sum_\ell c_\ell$ and is provably not slashed),
\decl{conservation} (every settled channel splits exactly $D$ between the
two parties, cooperative close and fund-slash forfeit alike), and
\decl{self\_slash\_race\_closed} (settlement happens at most once and no
path strands funds, so the self-slash race cannot leave the payee short).
The two deferrals the previous revision of this paper carried here are now
discharged. \texttt{Zkpc/Refund/Cascade.lean} models the full
failed-upgrade cascade --- successive withheld-receipt upgrades, one count
restored per round: upgrade claims never overshoot the certified count
(\decl{cascade\_upgrades\_le\_understatement}), a terminal cascade has
settled (\decl{cascade\_terminal\_settled}), settlement happens at exactly
the true count with exactly $n - j$ upgrades
(\decl{cascade\_settled\_upgrades\_eq}), and the final payouts conserve
$D$ (\decl{cascade\_final\_payouts}), with an executable driver
(\decl{execCascade\_progress}). \texttt{Zkpc/Refund/Fleet.lean} lifts the
single-channel results to interleaved multi-channel reachability and
aggregates no-overspend, settlement conservation, and the cooperative
payer floor across any finite fleet (\decl{fleet\_no\_overspend},
\decl{fleet\_conservation}, \decl{fleet\_payer\_floor}).
\texttt{Zkpc/Fleet/Recovery.lean} formalizes the fleet-side post-slash
recovery rule (MC19): pre-slash checkpoint eligibility
(\decl{preSlash\_claim\_eligible}, \decl{postSlash\_claim\_ineligible}),
sweep-before-conflict seniority, remainder-capped payouts
(\decl{identityRecovery\_capped}), exact conservation
(\decl{identityRecovery\_conservation}), full recovery when eligible
demand fits (\decl{identityRecovery\_all\_full}), and the distinct
fund-slash forfeit path (\decl{fundSlashRecovery\_full}).

\paragraph{Machine-checked wire-protocol bridges (the O1--O4 obligations,
discharged).} The gap the previous revision stated as an obligation --- T4
is proved on a proof-free view, the real wire carries a NIZK --- is now
closed for three concrete proof-bearing wire encodings, each with a
kernel-checked \emph{zero-loss} instance of \decl{zkBridgeObligation}
(\texttt{Zkpc/Games/FullTicketInstance.lean},
\texttt{Zkpc/Games/SigmaInstance.lean}): a \emph{masked-proof wire} (the
honest prover retains a private witness and emits it under a fresh
additive one-time mask; \decl{evalDist\_spendBatch\_maskedProof} proves
the witness-dependent real transcript equals the simulator distribution
exactly, giving \decl{T4\_maskedProof\_unlinkability} and
\decl{maskedProof\_zkBridge}); an \emph{interactive Sigma wire}
(\texttt{Zkpc/Crypto/LinearSigma.lean}: verifier completeness, a simulator
with exact transcript equality \decl{evalDist\_real\_eq\_simulated},
two-transcript special-soundness extraction \decl{special\_soundness};
\decl{T4\_sigmaFlat\_unlinkability}, \decl{sigmaFlat\_zkBridge}); and a
\emph{lazy-ROM Fiat--Shamir wire} (\texttt{Zkpc/Crypto/FSRom.lean}:
simulator distributions \decl{evalDist\_fsProveLazy\_eq\_simulated} plus
explicit programming- and fork-collision bounds
\decl{fsProgramCollisionBound}, \decl{fsForkChallengeCollisionBound};
\decl{T4\_fsFlat\_unlinkability}, \decl{fsFlat\_zkBridge}). The
refund-side B-instance obligations are likewise discharged
(\texttt{Zkpc/Games/BInstances.lean}): the rerandomized challenge path
(\decl{bRerand\_spendBatch\_none\_zero}, O2), adversary-issued genesis
receipts and issuer receipt updates with capability monotonicity
(\decl{bIdeal\_openCh\_adversary\_genesis},
\decl{bIdeal\_serve\_issuer\_receipt}, \decl{bIdeal\_serve\_capable\_mono},
M2/O3), and close-view simulatability for both B-static and B-rerand
(\decl{bIdeal\_closeViewSimulatable}, O4). On the refund cryptography
itself, \texttt{Zkpc/Crypto/MaskedEncryption.lean} proves exact
distributional rerandomization and refund-update privacy for the additive
masked cipher, and \texttt{Zkpc/Crypto/ReceiptMac.lean} proves a $1/|F|$
fresh-message forgery bound for the one-time algebraic receipt MAC
(\decl{mac\_forgery\_bound}). What remains beyond these
information-theoretic reference layers is deployment-grade: a concrete
hash-implementation reduction for the FS layer and multi-query EUF-CMA
signature/MAC reductions (\S\ref{sec:limits}).

\paragraph{Executable refinement.} The relational transition systems the
theorems quantify over are now refined by executable operations:
\texttt{Zkpc/Core/Refinement.lean} proves that executable open, honest
spend, fresh redeem, payer close, identity dispute, arbitrary sweep lists,
and the MC20 contract drivers (close dispute, successful settlement,
settlement-time voiding) return traces of their corresponding \decl{Step}
constructors (\decl{sweep\_refines\_trace},
\decl{refined\_steps\_reachable});
\texttt{Zkpc/Refund/Refinement.lean} and
\texttt{Zkpc/Fleet/Refinement.lean} do the same for refund
accept/close/force-close and fleet tick/admission/slash. States generated
by running the executable layer therefore inherit T1--T6 and the
refund/fleet invariants by construction.

\paragraph{The multi-recipient network layer (the named open problem,
definitional half).} \texttt{Zkpc/Network/State.lean} defines a
portable-deposit network --- one deposit funding arbitrarily many
recipients over a shared global nullifier set with recipient-directed
settlement --- and proves global deduplication (\decl{global\_dedup}),
network-wide no-overspend (\decl{no\_overspend}), exact
recipient-partitioned payout accounting, unrelated-recipient view
isolation, and executable refinement.
\texttt{Zkpc/Network/Credential.lean} gives the first concrete credential
adapter --- recipient, global nullifier, value, and payload bound into a
Fiat--Shamir statement --- with honest issuance verifying, verified fresh
redemption refining to network admission, cross-recipient nullifier replay
rejected (\decl{redeem\_rejects\_global\_replay}), and an end-to-end
payment theorem composing verification, executable redemption, settlement,
reachability, and shared-deposit no-overspend
(\decl{credential\_payment\_end\_to\_end}).
\texttt{Zkpc/Network/Issuance.lean} adds a finite threshold-issuance
reference: share aggregation correctness, perfectly hiding blind requests
(\decl{evalDist\_blindRequest\_uniform},
\decl{issuerView\_message\_independent}), fork extraction
(\decl{ticket\_fork\_extracts}), and exact recipient-view
simulation/unlinkability (\decl{recipientView\_unlinkable},
\decl{recipientView\_simulatable}).

\paragraph{One-trace composition.} \texttt{Zkpc/Core/Composition.lean}
bundles the guarantees that were previously separate endpoints.
\decl{channel\_endToEnd\_composition}: on a single reachable channel
trace, an honest payer with a posted close reaches a settled successor
state of the \emph{same} trace at which, simultaneously, the close settles
with the exact payer floor (T5+T3), every member satisfies no-overspend
(T1), the payee is settled exactly (T2), and the honest closer is
unslashed (exculpability). \decl{wire\_endToEnd\_composition}: for the
verified Fiat--Shamir wire family, perfect T4 unlinkability and the
zero-loss ZK bridge hold together. Both are axiom-clean.

\paragraph{Machine-checked nullifier-chain channel (instantiation C).}
The \S\ref{sec:constructions} design is formalized in
\texttt{Zkpc/Chain/} with the signature scheme idealized as transition
guards and the hash chain as a lazily-sampled random oracle
(collision-freedom carried as an explicit injectivity hypothesis on the
chain, stated where used). Safety (\texttt{State.lean}, Class A):
\decl{chain\_no\_overspend}, \decl{bob\_never\_loses} (honest close
pays exactly the closed balance --- \decl{honest\_close\_exact};
challenged stale close and Alice-AWOL timeout forfeit the whole deposit),
\decl{alice\_refund\_liveness} (a never-countersigned channel refunds
exactly $D$), \decl{conservation}, and \decl{no\_overpay\_recovery}.
The collision mechanism (\texttt{Collision.lean}) is proved in
\emph{both} directions: \decl{stale\_close\_detectable} (closing any
extended state --- the genesis-refund case uniformly included --- opens a
nullifier some message already revealed, so an honest Bob holds the
colliding challenge witness) and \decl{honest\_close\_unchallengeable}
/ \decl{honest\_close\_never\_slashed} (closing the latest
countersigned state opens a nullifier no message ever revealed --- the
design's exculpability, obtained without any FRAME-style probabilistic
argument), with the exactness lemma \decl{collision\_iff\_stale}
justifying the challenge guard. Per-request anonymity
(\texttt{Anonymity.lean}, Class B):
\decl{chain\_two\_payment\_anonymity} proves advantage exactly $0$
for every adversary in the two-payment linkage game
(same-chain-consecutive vs independent-channels), by coupling both worlds
to one canonical fresh view --- nullifiers as fresh-uniform oracle slots,
balance commitments as one-time additive masks. The game's docstring pins
what is \emph{not} covered: $\delta$-value correlation, timing, and the
base protocol's boundary leaks ($D$, close amounts, recipient,
footprint). An executable refinement (\texttt{Chain/Refinement.lean})
drives the machine.

An earlier review round on the game files produced a fix list --- the
adversary's view omitting the proof object (which would trivialize the
zero-knowledge step), the refund genesis receipt needing the adversary in
the loop, FRAME needing the close-time nullifier reveals --- and every item
on it was discharged before the proofs landed (\decl{zkBridgeObligation}
as the stated bridge, adversary-issued genesis, the \decl{nfAt} superset,
the \decl{roId} commitment surface).

What stands today: every one of the seven security statements is
machine-checked on the flat-ticket instantiation; the refund variant's
safety, conservation, failed-upgrade cascade, and finite-fleet aggregation
are machine-checked; the O1--O4 model-to-real bridges are discharged with
zero loss for three concrete wire encodings and the B instantiation; the
executable layer refines the relational one; the calibration battery
confirms the games are not vacuous; the channel and wire guarantees
compose on one trace; and the whole tree is axiom-clean (K2). The one
substantive in-model deferral left is stated where it occurs: T7's
unconditional bound awaits the secret-averaged handler coupling
(\decl{FrameDeferredSamplingAvg}), with the quantitative composition
endpoint already kernel-checked and the pointwise certificate variant
kernel-refuted. Everything else that remains is deployment-grade
cryptography --- concrete hash/signature reductions behind the ideal
reference layers --- and circuits are out of the model boundary
(\S\ref{sec:boundary}).

\section{Honest limits}\label{sec:limits}

\paragraph{Recipient-boundness.} The object binds a deposit to one logical
payee. That is the defining restriction, shared with ACT/ARC and every
keyed-verification design; the fleet setting relaxes it only by making $N$
gateways one logical payee with an internal consistency problem (which T6
prices). It is not a multi-merchant payment system.

\paragraph{Capital lockup.} $D$ per payer--payee pair, locked for the
channel's life. Below a payment-frequency threshold a channel does not pay
for its lockup~\cite{guasoni2024}; per-pair channels against $N$ verifiers
multiply this cost by $N$, which is the argument for the
fleet-as-one-payee shape.

\paragraph{Funding-graph leakage.} $\mathsf{Open}$ is a public ledger
event naming $cm$ and $D$; the deposit transaction links a funding
address --- and via exchange KYC trails, plausibly an identity --- to
membership. The sources prescribe shielded or Privacy-Pools-style funding;
we do not model it, and without it the anonymity set is fiction at the
funding edge. BOLT's anonymized-capital requirement is the same point one
decade earlier.

\paragraph{The spend-count-at-close side channel.} Both closes reveal the
spend count: as the certified $j$ in the refund variant, as
$\lfloor D/C \rfloor - |U|$ in the flat enumeration. T4 does not cover
it --- the game terminates at challenge delivery precisely because a
close after the challenge leaks the count. A member that closes
immediately after a distinctive burst correlates its count with observed
traffic. Mitigations (delay closes; close on round counts) are
operational, unproven here.

\paragraph{Within-epoch linkability, by design.} All of a member's spends
in one epoch share $nf_e$: rate limiting \emph{is} a linking mechanism at
epoch granularity. T4 claims unlinkability across epochs and to identity,
never within an epoch; the game's freshness condition is that scope made
formal. Epoch length is the privacy/throughput dial, and cross-epoch
intersection attacks over stable memberships remain --- as everywhere in
this literature --- an unpriced erosion.

\paragraph{A slash is retroactive deanonymization --- identity-slash only.}
An identity-slash publishes $k$, and from $k$ every $nf_i = H_{nf}(H_a(k,
i))$ and every epoch pseudonym $nf_e = H_e(k, e)$ is enumerable --- the
protocol relies on this for post-slash sweep attribution. So an
identity-slash does not cost the member $D$; it retroactively links the
member's \emph{entire lifetime request history} to its now-public $cm$. The
privacy T4 delivers is therefore contingent and revocable: it holds exactly
as long as $k$ stays secret, and the protocol itself contains the mechanism
that publishes $k$. There is no forward anonymity by design --- a legitimate
trade (it is what makes the slash claimable by anyone, watchtower-free), but
one a bare ``unlinkability'' claim would over-read. The blast radius is
widest for accidental double-emission: an honest-but-buggy client that
re-emits a different message at a used index self-slashes and loses both $D$
and its history; T7 says nothing here, because the client did double-sign.
Exculpability protects \emph{correct} users, not merely honest ones, and
this is the weight FRAME actually carries: it is the sole barrier between
every member and total retroactive deanonymization by the fleet's own
operators. The distinction of \S\ref{sec:object} bounds the damage --- a
fund-slash (false-claim void, settlement bar, B failed-upgrade forfeit)
keeps $k$ hidden and links nothing; only the evidence-pair identity-slash
deanonymizes. (Instantiation C is the design-space counterpoint: its
penalties are fund-forfeit only, so this entire limit is absent there at
the cost of interactivity and per-pair channels --- see
\S\ref{sec:constructions}.) Correspondingly, FRAME machine-checks the identity-slash door;
A's close-dispute exculpability is covered by the Core exculpability lemma
(\decl{honest\_never\_slashed}), while B's failed-upgrade exculpability
remains a specification-level argument, not machine-checked.

\paragraph{The stale-close residue in B.} Under receipt withholding, an
honest B payer can be wedged one count behind its true spend count, and its
only available close is on a stale receipt --- which the revealed $nf_j$
links to $cm$ for that one epoch session (rev-11 F10-1, MC15). The upgrade
sub-window restores the \emph{funds} but not the \emph{privacy}: the first
close is on-ledger and permanent. The nuance is that the linkage is not the
adversary's to take freely --- a payee declining to dispute keeps the
withheld receipt unpublished and gains the linkage at the cost of one
forgone $c \le C_{\max}$, so publication prices the payee's fund
\emph{recovery}, not the linkage; and the session extension of the linkage
requires the ticket transcript, which in B only the payee holds. It is one
session, not a lifetime, but it is a real residue and we state it.

\paragraph{From two candidates to the population, and the cost of aborts.}
T4 is a two-candidate game; deployments care about anonymity within the live
membership. The standard hybrid does go through here --- other members
simulate as corrupt payers run honestly, since payers hold independent
secrets, never interact peer-to-peer, and $D$ is a global constant --- so
the $2 \to n$ reduction loses the usual factor and no more; we state the
corollary but do not machine-check it, and note that with the
adversary-supplied genesis in B the hybrid must be checked to respect the
genesis stage (an unquantified erosion). The delivered guarantee is
indistinguishability \emph{within the challenge-capable set}, and that set
is adversary-controlled at zero protocol cost: refusal of service in A,
receipt withholding in B, each evicts a member from the set. After $q$
evictions the set is $n - q$, down to $1$, and nothing in T4 resists this
--- the defence is operational (members notice starvation and leave;
gateways compete), not cryptographic. No theorem in T1--T7 constrains a
payee's right to refuse service; the protocol prices refusal at zero.

\paragraph{Window recovery presumes fleet honesty.} Post-slash recovery
pays registered gateways against pre-slash checkpoints. A member colluding
with a corrupt gateway can pre-checkpoint real, cooperatively produced
``service'' at every solvent index and exhaust the slashed remainder
senior to honest gateways' claims. No theorem's adversary class covers
corrupt gateways in the \emph{recovery} role; T7 protects members from
gateways, not gateways from each other.

\paragraph{Close racing.} Service accepted between a close's ledger
inclusion and its settlement cannot be attributed to the closing channel:
in the refund variant the exposure is bounded by acceptance rate times
$\tau$ (pause acceptance while a close window is open); in the flat
variant an acceptance in flight at the close transaction is structurally
un-checkpointable and the tardy gateway bears exactly its un-checkpointed
tickets. Checkpoint cadence is each gateway's own recovery lever, and T2
is conditioned on it being exercised.

\paragraph{T6 is a bound, not a business case.} Recovery of the diverged
value is remainder-capped and checkpoint-gated; an exhaust-then-burst
member leaves little to recover. $f(L) < D$ bounds the burst by one
deposit and no more.

\paragraph{$x \neq 0$.} The signal at digest zero is the secret. The
deployment must domain-separate $H_x$ away from zero; the formalization
carries this as an explicit hypothesis, machine-checked in both directions
(\decl{rln\_single\_point\_hiding} requires it;
\decl{rln\_x\_zero\_degenerate} shows why).

\paragraph{The named open problem: multi-recipient generalization.}
Everything in the channel object binds value to one logical payee. A
member paying $N$ \emph{independent} payees with per-pair channels needs
$N$ deposits, $N$ partitioned anonymity sets, and hands each payee BOLT's
abort lever --- the three costs that made the per-pair shape lose to the
fleet shape in our application analysis. The generalization needs either
portable deposits (value that moves between payees without a linking
event --- the hub line's problem, with the hub's anonymity accounting) or
threshold issuance over a shared spent structure (the Nym-shaped
hybrid~\cite{nymzknym}, whose reconciliation guarantees are exactly what
would need the T6 treatment). PrivateX402~\cite{privatex402} shows what
giving up costs: one deposit across $N$ recipients, and every recipient
links every request. The definitional and accounting half of this problem
is now formalized (\S\ref{sec:results}): the portable-deposit network
machine, its credential adapter, and the threshold-issuance reference
construction prove global deduplication, no-overspend, payout
partitioning, recipient-view isolation, blind-request hiding, and
recipient-view unlinkability. What we still do not claim is the
composition that would make it a \emph{solved} problem: an adaptive
multi-session network game connecting those per-session distributions to
the executable admission and settlement trace, and a production
threshold-signature unforgeability reduction. The problem is now
half-closed, and we say which half.

Also outside every theorem: static corruption only; no clock skew (a
skew-tolerant model enlarges the T6 budget by a small constant); and the
entire traffic-analysis surface --- timing, size, content fingerprints ---
which in the source thread's own discussion is ``a richer fingerprint than
the wallet address itself.''

\section{Reproducibility}\label{sec:repro}

Repository: \url{https://github.com/dmarzzz/zk-payments-confetti}.
Toolchain, pinned: \texttt{leanprover/lean4:v4.30.0}; mathlib
\texttt{v4.30.0} (manifest revision \texttt{c5ea0035\ldots}); VCV-io at
commit \texttt{8f5dc4f2\ldots}~\cite{vcvio2026}. Build:
\begin{center}
\texttt{lake exe cache get \&\& lake build}
\end{center}
CI runs the build on the pinned toolchain and additionally fails on:
\texttt{sorry} anywhere; \texttt{axiom} outside
\texttt{Zkpc/Assumptions.lean}; \texttt{admit} or \texttt{native\_decide}
anywhere.

\begin{footnotesize}
\begin{longtable}{>{\raggedright\arraybackslash}p{0.17\textwidth}>{\raggedright\arraybackslash}p{0.24\textwidth}>{\raggedright\arraybackslash}p{0.51\textwidth}}
\caption{Theorem-to-file map.}\label{tab:map}\\
\toprule
Statement & File & Declarations \\
\midrule
\endfirsthead
\toprule
Statement & File & Declarations (cont.) \\
\midrule
\endhead
\bottomrule
\endfoot
T1 & \texttt{Zkpc/Core/T1.lean} & \decl{T1\_no\_overspend}, \decl{reach\_inv} \\
Exculpability (symbolic) & \texttt{Zkpc/Core/T1.lean} & \decl{honest\_never\_slashed} \\
T2 & \texttt{Zkpc/Core/T2.lean} & \decl{T2\_upper}, \decl{T2\_paid\_exact}, \decl{T2\_swept\_accepted}, \decl{sweepOne\_enabled}, \decl{T2\_collectable}, \decl{T2\_settles\_exactly} \\
T3 & \texttt{Zkpc/Core/T3.lean} & \decl{payer\_pay\_inv}, \decl{settleClose\_enabled}, \decl{T3\_settled\_amount}, \decl{T3\_payer\_balance\_security} \\
T5 & \texttt{Zkpc/Core/T5.lean} & \decl{T5\_payer\_close\_liveness}, \decl{settleClose\_stable}, \decl{tick\_progress} \\
T6 & \texttt{Zkpc/Fleet/T6.lean}, \texttt{Zkpc/Fleet/Basic.lean} & \decl{T6\_priced\_divergence}, \decl{T6\_accept\_count}, \decl{T6\_slash\_within\_L}, \decl{card\_le\_solvency\_of\_conflictFree}, \decl{card\_le\_rate\_window}, \decl{epochs\_in\_window}, \decl{fleet\_inv} \\
RLN algebra & \texttt{Zkpc/Games/RLN.lean} & \decl{rln\_recover\_a}, \decl{rln\_recover\_k}, \decl{rln\_single\_point\_hiding}, \decl{rln\_x\_zero\_degenerate}, \decl{rln\_evidence\_complete}, \decl{rln\_evidence\_sound} \\
Game framework & \texttt{Zkpc/Games/Framework.lean} & \decl{guessGap}, \decl{guessGap\_eq}, \decl{hiddenBitAdvantage\_eq\_half\_boolDistAdvantage}, \decl{hiddenBitAdvantage\_const}, \decl{hiddenBitAdvantage\_eq\_zero\_of\_distEquiv}, \decl{ChalAdversary}, \decl{withEvict} \\
T4 (unlinkability) & \texttt{Zkpc/Games/T4.lean}, \texttt{Zkpc/Games/Unlink.lean} & \decl{T4\_flat\_unlinkability} ($= 0$), \decl{unlinkGame}, \decl{unlinkAdvantage}, \decl{UnlinkScheme}, \decl{flat\_closeViewSimulatable} \\
T7 (framing) & \texttt{Zkpc/Games/T7.lean}, \texttt{Zkpc/Games/Frame.lean} & \decl{T7\_frame\_bound} ($\le 1/|F|$ under \decl{hobliv}), \decl{frameWinProb\_YK\_eq\_one}, \decl{frameWinProb\_aReuse\_eq\_one}, \decl{frameWinProb\_slopeReveal\_eq\_one}, \decl{frameGame}, \decl{frameWinProb}, \decl{Slashes} \\
T7 query-budget composition & \texttt{Zkpc/Games/T7.lean}, \texttt{Zkpc/Games/Frame\{Audit,Ideal,Deferred\}.lean} & \decl{FrameQueryBounds}, \decl{uniformSecretProbeBound}, \decl{uniformSlopeProbeBound}, \decl{T7\_frame\_bound\_of\_pointwise}, \decl{T7\_frame\_query\_bound}, \decl{frameDeferredSampling\_refuted}, \decl{FrameDeferredSamplingAvg}, \decl{T7\_frame\_query\_bound\_avg}, \decl{auditedFrameImpl}, \decl{idealFrameImpl}, \decl{FrameCoupled}, \decl{idealizeFrame} \\
Wire ZK bridges (O1) & \texttt{Zkpc/Games/FullTicketInstance.lean}, \texttt{Zkpc/Games/SigmaInstance.lean} & \decl{T4\_maskedProof\_unlinkability}, \decl{maskedProof\_zkBridge}, \decl{T4\_sigmaFlat\_unlinkability}, \decl{sigmaFlat\_zkBridge}, \decl{T4\_fsFlat\_unlinkability}, \decl{fsFlat\_zkBridge}, \decl{fullFlat\_zkBridge} \\
Sigma / Fiat--Shamir cores & \texttt{Zkpc/Crypto/LinearSigma.lean}, \texttt{Zkpc/Crypto/FSRom.lean} & \decl{completeness}, \decl{evalDist\_real\_eq\_simulated}, \decl{special\_soundness}, \decl{evalDist\_fsProveLazy\_eq\_simulated}, \decl{fsProgramCollisionBound}, \decl{fsForkChallengeCollisionBound} \\
B-instance obligations (O2--O4) & \texttt{Zkpc/Games/BInstances.lean} & \decl{bRerand\_spendBatch\_none\_zero}, \decl{bIdeal\_openCh\_adversary\_genesis}, \decl{bIdeal\_serve\_issuer\_receipt}, \decl{bIdeal\_serve\_capable\_mono}, \decl{bIdeal\_closeViewSimulatable} \\
Refund crypto references & \texttt{Zkpc/Crypto/MaskedEncryption.lean}, \texttt{Zkpc/Crypto/ReceiptMac.lean} & \decl{evalDist\_rerandomize\_cipher\_uniform}, \decl{evalDist\_refundUpdate\_cipher\_uniform}, \decl{mac\_forgery\_bound} \\
Refund cascade & \texttt{Zkpc/Refund/Cascade.lean} & \decl{cascade\_upgrades\_le\_understatement}, \decl{cascade\_settled\_upgrades\_eq}, \decl{cascade\_terminal\_settled}, \decl{cascade\_final\_payouts}, \decl{execCascade\_progress} \\
Refund fleet aggregation & \texttt{Zkpc/Refund/Fleet.lean} & \decl{fleet\_no\_overspend}, \decl{fleet\_conservation}, \decl{fleet\_payer\_floor} \\
Fleet recovery (MC19) & \texttt{Zkpc/Fleet/Recovery.lean} & \decl{identityRecovery\_conservation}, \decl{identityRecovery\_capped}, \decl{identityRecovery\_all\_full}, \decl{preSlash\_claim\_eligible}, \decl{postSlash\_claim\_ineligible}, \decl{fundSlashRecovery\_full} \\
Executable refinement & \texttt{Zkpc/\{Core,Refund,Fleet\}/Refinement.lean} & \decl{sweep\_refines\_trace}, \decl{refined\_steps\_reachable}, \decl{exec*\_refines\_step}, \decl{exec\_step\_reachable} \\
Multi-recipient network & \texttt{Zkpc/Network/\{State,Credential,Issuance\}.lean} & \decl{no\_overspend}, \decl{global\_dedup}, \decl{redeem\_rejects\_global\_replay}, \decl{credential\_payment\_end\_to\_end}, \decl{evalDist\_blindRequest\_uniform}, \decl{ticket\_fork\_extracts}, \decl{recipientView\_unlinkable} \\
One-trace composition & \texttt{Zkpc/Core/Composition.lean} & \decl{channel\_endToEnd\_composition}, \decl{wire\_endToEnd\_composition} \\
Nullifier-chain channel (C) & \texttt{Zkpc/Chain/\{State,Collision,Anonymity,Refinement\}.lean} & \decl{chain\_no\_overspend}, \decl{bob\_never\_loses}, \decl{honest\_close\_exact}, \decl{alice\_refund\_liveness}, \decl{conservation}, \decl{no\_overpay\_recovery}, \decl{stale\_close\_detectable}, \decl{honest\_close\_unchallengeable}, \decl{collision\_iff\_stale}, \decl{chain\_two\_payment\_anonymity} \\
Calibration & \texttt{Zkpc/Games/Calibration.lean} & \decl{unlinkAdvantage\_staticDistinguisher\_eq\_half}, \decl{unlinkAdvantage\_bRerand\_eq\_zero}, \decl{unlinkAdvantage\_aIndexLeak}, \decl{unlinkAdvantage\_nfeReuse}, \decl{unlinkAdvantage\_multTagDistinguisher\_eq\_half} \\
Refund variant (B, $N{=}1$) & \texttt{Zkpc/Refund/Safety.lean}, \texttt{Zkpc/Refund/State.lean} & \decl{T1\_B\_no\_overspend}, \decl{T3\_B\_floor}, \decl{conservation}, \decl{self\_slash\_race\_closed} \\
Assumption registry & \texttt{Zkpc/Assumptions.lean} & \decl{Named}, \decl{dischargedBy} (no \texttt{axiom} declarations exist) \\
State machines & \texttt{Zkpc/Core/State.lean}, \texttt{Zkpc/Core/Flat.lean}, \texttt{Zkpc/Fleet/Basic.lean} & transition systems the above quantify over \\
\end{longtable}
\end{footnotesize}

The specification of record is \texttt{Spec.md} (revision 11); every Lean
definition is traceable to it, and the full adversarial review record ---
eleven rounds against the specification, three against the Lean games, plus
independent statement (K1), axiom (K2), and external-cryptographer (K4)
audits, with every counterexample --- is
\texttt{research\_knowledge/gates.md}. TLA+
models of the flat and fleet state machines, including ablation
configurations that replay the gateway-binding and merge-evidence
counterexamples (\texttt{tla/ZkpcFleetNoBind.cfg},
\texttt{tla/ZkpcFleetNoMergeEv.cfg}), are in \texttt{tla/}.

\bibliographystyle{plain}
\bibliography{refs}

\end{document}
